% META: {"id":"1SPE-SECDEG-EX-009","chapitre":"1SPE-SECOND-DEGRE","type_objet":"exercice","capacites":["1SPE-SECOND-DEGRE-C3"],"capacites_codes":["C3"],"methodes":["M3"],"parcours":1,"competences":["calculer","raisonner"],"duree_min":12,"mode_creation":"ex_nihilo","sources_inspiration":[],"parametres_sympy":{},"fichier_tex":"chapitres/1SPE-SECOND-DEGRE/exercices/1SPE-SECDEG-EX-009.tex","corrige_tex":"chapitres/1SPE-SECOND-DEGRE/corriges/1SPE-SECDEG-CO-009.tex","status":"approved","origin":"generated","status_history":["generated","audited","accepted","approved"],"release_acceptance":"RELEASE_OWNER_BATCH_ACCEPTANCE","acceptance_closure_digest":"sha256:9b3ccf9a81c5520fbb7e03b7b2d3e7bf2057a02834b6d908bf3be5f49f3b553f","acceptance_receipt_digest":"sha256:4fad86f0282e0fa4e0419cca1ad6379ae7af5a280c04a4a90880f90a1c044242"}
% BEGIN-VERIFY
% from sympy import symbols, solve, Rational, sqrt, simplify
% x = symbols('x')
% # Equation ramenee a la forme ax^2+bx+c=0
% # (x+1)(x-3) = -3  =>  x^2 - 2x - 3 + 3 = 0  =>  x^2 - 2x = 0
% f = x**2 - 2*x
% sols = solve(f, x)
% assert sorted(sols) == [0, 2]
% # Verification
% assert (0+1)*(0-3) == -3
% assert (2+1)*(2-3) == -3
% # Equation 2x^2 - 5x = 3  =>  2x^2 - 5x - 3 = 0
% g = 2*x**2 - 5*x - 3
% delta_g = 25 + 24
% assert delta_g == 49
% sols_g = solve(g, x)
% assert sorted(sols_g) == [Rational(-1,2), 3]
% END-VERIFY
\begin{exercice}{1SPE-SECDEG-EX-009}{1}{12}
Résoudre dans $\mathbb{R}$ les équations suivantes. On ramènera chacune à la forme $ax^2 + bx + c = 0$ avant de calculer $\Delta$.

\begin{enumerate}
  \item $(x+1)(x-3) = -3$
  \item $2x^2 - 5x = 3$
\end{enumerate}
\end{exercice}
